% ===================================================================
% CHAPTER 2: PROBLEM FORMULATION
% ===================================================================
\chapter{Формулација на проблемот}
\label{chap:problem}

Федерираното учење го решава проблемот со приватноста при обуката на модели за откривање
упади, но воведува нов проблем при нивното \textbf{одржување}. Кога еден јазол ќе биде
идентификуван како компромитиран, или кога сопственикот на податоците ќе го повлече
својот пристанок, влијанието на тие податоци мора да биде отстрането од моделот. Бришењето
на суровите записи не е доволно --- моделот ги задржува во своите тежини информациите
научени од нив. Ова поглавје го формализира проблемот и ги поставува истражувачките
прашања што го водат остатокот од трудот.

\section{Формална поставка}

Нека федерацијата се состои од $K$ клиенти, при што клиентот $k$ располага со локално
податочно множество $D_k = \{(x_i, y_i)\}_{i=1}^{n_k}$, каде $x_i \in \mathbb{R}^d$ е
вектор на мрежни обележја, а $y_i \in \{0, 1\}$ ја означува етикетата \emph{бенигно} наспроти
\emph{напад}. Глобалниот модел со параметри $w$ се добива со алгоритмот FedAvg
\cite{mcmahan2017fedavg}: во рундата $t$ секој клиент ги презема тековните глобални
параметри $w^{(t)}$, изведува локална оптимизација врз $D_k$ и враќа
$w_k^{(t+1)}$, а серверот пресметува тежинска средна вредност

\[
w^{(t+1)} \;=\; \sum_{k=1}^{K} \frac{n_k}{\sum_{j} n_j}\, w_k^{(t+1)} .
\]

Во моментот $\tau$ се утврдува дека подмножеството $D_f \subset D_k$ на некој клиент $k$ е
компромитирано. Барањето е да се произведе модел што е \textbf{ослободен од влијанието} на
$D_f$. Нека $\mathcal{A}$ е постапката за обука и $\mathcal{U}$ постапката за заборавање.
Заборавањето е \textbf{егзактно} ако распределбата на параметрите произведени од
$\mathcal{U}$ врз моделот обучен на $D_k$ е идентична со распределбата што би ја дала
повторна обука од нула само врз задржаните податоци:

\[
\mathcal{U}\bigl(\mathcal{A}(D_k),\, D_f\bigr) \;\stackrel{d}{=}\; \mathcal{A}\bigl(D_k \setminus D_f\bigr).
\]

Референтното решение што тривијално ја задоволува оваа равенка е \textbf{наивното повторно
тренирање} --- десната страна се пресметува буквално. Алтернативата, \textbf{SISA}
\cite{bourtoule2021sisa}, го дели локалното множество на $S$ изолирани шардови, секој
поделен на $R$ сегменти, и ја чува состојбата по секој сегмент; при барање за заборавање
се враќа наназад само засегнатиот шард и се обработуваат повторно само сегментите по
точката на компромитирање.

\subsection{Опсег на гаранцијата}

Клучно ограничување што мора да се наведе однапред: и двата пристапа обезбедуваат
\textbf{егзактно локално} заборавање --- обновената состојба на клиентот е докажливо ослободена
од влијанието на $D_f$ --- но само \textbf{приближно глобално} заборавање. FedAvg е агрегација
без состојба; откако придонесот на клиентот $k$ е веќе вграден во $w^{(t)}$, ниту една
постапка изведена на страната на клиентот не може да го отстрани од глобалниот модел.
Оттука, споредбата во овој труд \textbf{не е} „иста работа со различен трошок“, туку
\textbf{трошок за постигнување на иста гаранција}.

\section{Цената на заборавањето како предмет на мерење}

Постојната литература за заборавање ја изразува цената во пресметковни единици --- епохи,
рунди, операции со подвижна запирка --- и ја мери во симулација. На вистински уред на работ
таа апстракција се распаѓа. Уредот од класата Raspberry Pi е ограничен истовремено по
пресметка, топлина, влезно-излезна пропусност и напојување, и тие ограничувања меѓусебно
се засилуваат: продолжена пресметка го загрева чипот, загревањето предизвикува термичко
придушување (\emph{thermal throttling}), а придушувањето ја продолжува пресметката.

Затоа овој труд ја дефинира \textbf{цената на заборавањето} не како скалар, туку како вектор
од физички величини измерени врз самиот хардвер:

\[
\mathbf{c} \;=\; \bigl(\underbrace{t_{\text{обн}}}_{\text{време}},\;
\underbrace{E_{\text{нето}}}_{\text{енергија}},\;
\underbrace{s_{\text{прид}}}_{\text{топлина}},\;
\underbrace{B_{\text{SD}}}_{\text{I/O}}\bigr),
\]

дополнет со \textbf{режискиот трошок во фазата на обука} што SISA го плаќа непрекинато
(запишување контролни точки во секоја рунда), и со \textbf{корисноста} на обновениот модел.
Вкупниот трошок на сопственост е збир од режискиот трошок и трошокот на обновувањето --- не
само второто.

\subsection{Рамка: трошок на отстранувањето, не сила на нападот}

Труењето со превртување на етикети во овој труд има точно определена улога: тоа
\textbf{дефинира јасно специфицирано множество примери што мора да се заборават}, и ништо
повеќе. Отстранувањето е мотивирано од фактот дека податоците се \textbf{идентификувани како
компромитирани} --- преку откривање или преку регулаторно право на заборавање
\cite{gdpr2016} --- а не од нивната измерена штета врз моделот. При скалата на овој
експеримент ограниченото труење нема мерлив ефект врз глобалната корисност, бидејќи FedAvg
го разредува придонесот на единствениот компромитиран клиент. Тоа не е слабост на
дизајнот, туку потврда на рамката: прашањето не е \emph{колку штети нападот}, туку
\textbf{колку чини отстранувањето}. Отпорноста на федерацијата кон просторно ограничено труење
е самостоен наод и се известува како таков.

Откривањето на компромитираните примери е \textbf{надвор од опсегот} на трудот; се користи
оракул-маска врз нив.

\section{Цели и истражувачки прашања}

Главното истражувачко прашање гласи:

\begin{quote}
\textbf{Колкави се физичките трошоци на ресурсите --- време, енергија, топлина и влезно-излезни
операции --- на заборавањето засновано на SISA наспроти наивното повторно тренирање, кога
ресурсно ограничен федериран јазол на работ мора да отстрани идентификувано
компромитирани податоци?}
\end{quote}

Тоа прашање се разложува на четири потпрашања, секое поврзано со претходно регистрирани
хипотези и со мерливи зависни променливи:

\begin{itemize}
  \item \textbf{RQ1 --- Пресметковен и енергетски трошок.} Колку се разликуваат наивното повторно
  тренирање и SISA по \textbf{време до обновување} и по \textbf{нето потрошена енергија} (со
  одземена основна потрошувачка во мирување) на реален, термички ограничен уред на работ?
  \emph{(Хипотези H1 и H2.)}
  \item \textbf{RQ2 --- Термичко и влезно-излезно однесување.} Како двата пристапа го оптоваруваат
  \textbf{термичкиот буџет} на уред без активно ладење, и како го прераспределуваат трошокот
  кон \textbf{влезно-излезните операции} на SD-картичката? \emph{(Хипотези H3 и H4.)} Ова
  прашање бара внимателен избор на показател: врвната температура се заситува за секое
  одржливо оптоварување и затоа не разликува, па термичкиот сигнал го носи \textbf{траењето}
  на придушувањето.
  \item \textbf{RQ3 --- Задржување на корисноста.} Дали обновениот глобален модел ја задржува
  \textbf{способноста за откривање}, мерена со показатели отпорни на класна нерамнотежа ---
  специфичност, избалансирана точност и Метјуов коефициент на корелација --- наместо со
  одѕив и F1, кои се тривијални врз множество со 96,5\,\% напади \cite{chicco2020mcc,
  saito2015prc}? \emph{(Хипотеза H5.)}
  \item \textbf{RQ4 --- Вкупен трошок на сопственост.} Каков е нето биланс на SISA кога режискиот
  трошок за запишување контролни точки во текот на \textbf{целата фаза на обука} ќе се додаде
  на заштедата остварена при обновувањето? \emph{(Хипотеза H6.)}
\end{itemize}

Целта на трудот не е да се утврди дека еден од двата пристапа е подобар во секој поглед,
туку да се \textbf{измери компромисот} меѓу нив на вистински хардвер, и тоа да се направи со
доволна статистичка строгост за наодот да биде употреблив при проектирање на производни
федерирани системи на работ на мрежата.
